% lint-skip: 未完成ステータス（figure/table caption 未対応，人手チェック時に解消予定）
\documentclass[../../main.tex]{subfiles}
\begin{document}

\begin{flushright}
  \footnotesize\textit{【自動文字起こし・要確認】}
\end{flushright}

[解] $\dfrac{1}{2}$ の確率で1点増し，$\dfrac{1}{2}$ の確率で1点減る．

$n$ 回目までに奇数の目が出た回数 $X(n)$ ，$n$ 回目までに偶数の目が出た回数 $Y(n)$ とおくと $n \le 5$ の時，常に
\[ X(n)+2 > Y(n) \quad (\because X(n)+Y(n)=n) \]
がみたされていれば良い ($n=5$ では等号が成立しても良い)．このような場合の数は，下の図で，格子の上を通って $A$ から各終点まで行く最短経路のうち $Y=X+2$ を通らないものと1対1対応している．

\begin{figure}[htb]
  \centering
  \begin{tikzpicture}[x=0.85cm, y=0.85cm, >=stealth,
    every node/.style={font=\small}]
    \draw[step=1, very thin, gray!45] (0,0) grid (5,5);
    \draw[->, thick, gray!80!black] (-0.2,0) -- (5.8,0)
    node[right] {$X$};
    \draw[->, thick, gray!80!black] (0,-0.2) -- (0,5.7)
    node[above] {$Y$};
    \draw[line width=1.3pt, red!80!black] (0,2) -- (3,5);
    \node[above right=2pt, red!80!black] at (2.4,4.4)
    {$Y=X+2$};
    \fill[blue!85!black] (0,0) circle (1.5pt);
    \node[below left=2pt] at (0,0) {$A$};
    \foreach \x/\y/\score in {5/0/7,4/1/5,3/2/3,2/3/1} {
      \fill[teal!75!black] (\x,\y) circle (2pt);
      \node[above right=2pt, teal!75!black] at (\x,\y)
      {\score 点};
    }
  \end{tikzpicture}
  \caption{境界 $Y=X+2$ を通らずに達する5回後の4つの得点}
\end{figure}

7点 $\cdots$ 1通り

5点 $\cdots$ 5通り

3点 $\cdots$ 9通り

1点 $\cdots$ 5通り

だから，
5回振ることのできる確率は $\dfrac{20}{2^5} = \dfrac{5}{8}$ 終

得点の期待値は $\dfrac{1}{2^5} (7+25+27+5) = \underline{2 \text{点}}$ 終

\end{document}
